\documentclass[10pt,a4paper]{amsart}
\usepackage[margin=19mm]{geometry}
\usepackage{amsmath,amssymb,mathtools}
\usepackage[scr=rsfs,cal=euler]{mathalfa}
\usepackage{xfrac}
\usepackage[unicode,hidelinks,bookmarks=false]{hyperref}
\newcommand{\ie}{i.e.\ }
\newcommand{\cf}{cf.\ }
\newcommand{\resp}{respectively\ }
\newcommand{\Subsection}{\subsection}
\providecommand{\nobreakdash}{\nobreak\hbox{-}}
\setlength{\emergencystretch}{2em}
\multlinegap0pt
\usepackage{amssymb}
\usepackage{xcolor}
\newtheorem{theorem}{Theorem}
\title{A modified Brinkman penalization fictitious domain method for the unsteady Navier-Stokes equations}
\author{Zhanybek Baitulenov, Maxim Olshanskii, Almas Temirbekov, Nurlan Temirbekov, Syrym Kasenov}
\date{Source: arXiv:2512.19580v1 (CC BY 4.0)}
\begin{document}
\maketitle

\noindent\emph{Source and licence.} A modified Brinkman penalization fictitious domain method for the unsteady Navier-Stokes equations. Version arXiv:2512.19580v1 (CC BY 4.0), distributed by arXiv under the Creative Commons Attribution 4.0 International licence, http://creativecommons.org/licenses/by/4.0/, as recorded in the arXiv licence metadata for this exact version and shown on the abstract page https://arxiv.org/abs/2512.19580v1. Full licence notice: \texttt{licenses/arxiv-2512.19580-CC-BY-4.0.txt}.\par\smallskip

\noindent\textit{Editorial scope.} Counted unit: the article's theorem 15 (source0.tex lines 304--314). The article's complete original proof of it is delivered with it (source0.tex lines 316--470, one \texttt{proof} environment of 155 lines). No mathematics is written, repaired or re-derived: the statement and the proof are the article's own lines, in the article's own notation, and no retained line is substituted or re-flowed.  Retained uncounted as the source-local context the proof needs: lines 62--64; lines 69--71; lines 107--109; lines 113--115.  Nothing was omitted from any retained span, so the package carries no omission record.  The retained lines cite 1 works, whose entries are transcribed verbatim from the article's own bibliography source in the e-print and delivered with short editorial labels recorded in the item's reference mapping.  No retained line contains a figure environment, an image-inclusion command or drawing code, so no image is delivered and the package declares no asset dependency.  Editorial formatting only, in addition: an AMS \texttt{amsart} preamble that restates the article's own declarations and macros used by the retained lines, namely the environments \texttt{theorem} on separate counters, together with the standard packages those lines need.  The retained environments print in this document as Theorem 1 in order of appearance, whereas the article prints the same items with the numbering of its own full document.  The complete native source of the article as distributed in the arXiv e-print for this exact version is bundled unchanged as \texttt{sources/191489/source0.tex}.
\begin{equation} \label{St 1}
	\frac{\partial v}{\partial t} + (v\cdot\nabla)v = \mu\Delta v - \nabla P + f,
\end{equation}

\begin{equation} \label{St 3}
	v|_{t=0} = v_0(x),\qquad v|_{S} = 0.
\end{equation}

\begin{equation} \label{St 8}
	\frac{\partial v^\varepsilon}{\partial t} + (v^\varepsilon \cdot \nabla)v^\varepsilon = \mu\Delta v^\varepsilon - \nabla P^\varepsilon - \frac{\xi(x)v^\varepsilon}{\varepsilon\|v^\varepsilon\|_{L_2(D_1)}^{\beta}} + f^{\varepsilon},
\end{equation}

\begin{equation} \label{St 10}
	v^\varepsilon|_{t=0} = v_0(x),\qquad v^\varepsilon|_{S_1} = 0.
\end{equation}

\begin{theorem} Assume $f \in L_2(0, T; L_2(\Omega))$ and $v_0(x) \in V(\Omega)$, then for  a sufficiently small  $ T_0>0$, there exists the unique solution $ v^{\varepsilon}$ to \eqref{St 8}--\eqref{St 10} such that 
	$ v^{\varepsilon}\in L_2 (0,T_0; W^{2}_{2} (D))\cap L_\infty (0,T_0; V(D))$, $v^{\varepsilon}_{t}\in L_2 (0,T_0 ; L_2 (D))$ and   the following estimate holds: 
\begin{multline} \label{St 15}
\| v^{\varepsilon} \|^{2}_{L_\infty (0,T_0; V(D))} + \frac{1}{\varepsilon} \|v^{\varepsilon} \|^{2 - \beta}_{L_\infty (0,T_0; L_2 (D_1))} + \\
+ \|v^{\varepsilon}_{t} \|^{2}_{L_2 (0,T_0 ; L_2 (D))} 
+ \| v^{\varepsilon} \|^{2}_{L_2 (0,T_0; W^{2}_{2} (D))} 
+ \frac{1}{\varepsilon} \int_{0}^{T_0} \|\nabla v^{\varepsilon} \|^{2-\beta}_{L_2 (D_1)} \, dt 
\leq C,
\end{multline}
{with some $C$ depending on $f$, $\mu$, $\Omega$, and $D$.}
\end{theorem}

\begin{proof} % --- assumptions: smooth domain, div-free fields, homogeneous BCs for Stokes operator ---
	Denote by $\tilde\Delta $ the Stokes operator.  
Since $D$ has a smooth boundary, by the Stokes regularity we have
	\[
	\|v\|_{W^2_2(D)\cap V(D)} \le C_{\rm reg}\,\|\tilde\Delta v\|_{L_2(D)}
	\]
for $v\in W^2_2(D)\cap V(D)$.
	Hence, using interpolation and Sobolev embeddings,
	\[
	\|\nabla v\|_{L^3(D)} \le \|\nabla v\|_{L^2(D)}^{1/2}\|\nabla v\|_{L^6(D)}^{1/2}
	\le C\,\|\nabla v\|_{L^2(D)}^{1/2}\|v\|_{W^2_2(D)}^{1/2}
	\le C\,\|\tilde\Delta v\|_{L^2(D)}^{1/2}\|\nabla v\|_{L^2(D)}^{1/2}.
	\]
	
	Taking the $L_2(D)$-inner product of \eqref{St 8} with $v_t$  yields
	\[
	\|v_t\|_{L_2(D)}^2 - \frac{\mu}{2}\frac{d}{dt}\|\nabla v\|_{L_2(D)}^2
	+ \int_D (v\cdot\nabla)v\cdot v_t\,dx
	= \int_D f\,v_t\,dx - \frac{1}{\varepsilon\|v\|_{L_2(D_1)}^\beta}\int_{D_1} v\cdot v_t\,dx.
	\]
We estimate the nonlinear term by the Gagliardo--Nirenberg and Young inequalities:
	\[
	\begin{aligned}
		\Big|\int_D (v\cdot\nabla)v\cdot v_t\,dx\Big|
		&\le \|v_t\|_{L_2(D)}\|v\|_{L_6(D)}\|\nabla v\|_{L_3(D)} \\
		&\le C\|v_t\|_{L_2(D)}\|\nabla v\|_{L_2(D)}\|\tilde\Delta v\|_{L_2(D)}^{1/2}\|\nabla v\|_{L_2(D)}^{1/2} \\
		&\le \delta_1\|v_t\|_{L_2(D)}^2 + C_{\delta_1}\|\tilde\Delta v\|_{L_2(D)}\|\nabla v\|_{L_2(D)}^3 \\
		&\le \delta_1\|v_t\|_{L_2(D)}^2 + \theta_1\|\tilde\Delta v\|_{L_2(D)}^2 + C_{\theta_1,{\delta_1}}\|\nabla v\|_{L_2(D)}^6.
	\end{aligned}
	\]
	Similarly,
	\[
	\int_D f v_t \,dx \le \delta_2\|v_t\|_{L_2(D)}^2 + C_{\delta_2}\|f\|_{L_2(D)}^2.
	\]
	The damping-derivative term is exactly
	\[
	\frac{1}{\varepsilon\|v\|_{L_2(D_1)}^\beta}\int_{D_1} v\cdot v_t\,dx
	= \frac{1}{2\varepsilon(2-\beta)}\frac{d}{dt}\|v\|_{L_2(D_1)}^{2-\beta}.
	\]
	Choosing $\delta_1+\delta_2<1$ we obtain
	\begin{equation}\label{St 17}
		\|v_t\|_{L_2(D)}^2 + \frac{d}{dt}\Big({\frac\mu2}\|\nabla v\|_{L_2(D)}^2 + \frac{1}{\varepsilon}\|v\|_{L_2(D_1)}^{2-\beta}\Big)
		\le \theta_1\|\tilde\Delta v\|_{L_2(D)}^2 + C\big(\|\nabla v\|_{L_2(D)}^6 + \|f\|_{L_2(D)}^2\big).
	\end{equation}
	Next we take the $L_2(D)$-inner product of \eqref{St 8} with $\tilde\Delta v$.
 Then
	\[
	{\mu}\|\tilde\Delta v\|_{L_2(D)}^2 - \frac{1}{\varepsilon\|v\|_{L_2(D_1)}^\beta}\int_{D_1} v\cdot\tilde\Delta v\,dx
	= \int_D (v_t+(v\cdot\nabla)v-f)\cdot\tilde\Delta v\,dx.
	\]
	Estimating the right-hand side by Young's inequality gets:
	\[
	\begin{aligned}
		\int_D v_t\cdot\tilde\Delta v &\le \theta_2\|\tilde\Delta v\|_{L_2}^2 + C_{\theta_2}\|v_t\|_{L_2}^2,\\
		\int_D (v\cdot\nabla)v\cdot\tilde\Delta v &\le \theta_3\|\tilde\Delta v\|_{L_2}^2 + C_{\theta_3}\|\nabla v\|_{L_2}^6,\\
		\int_D f\cdot\tilde\Delta v &\le \theta_4\|\tilde\Delta v\|_{L_2}^2 + C_{\theta_4}\|f\|_{L_2}^2.
	\end{aligned}
	\]
	For the damping term on the left, we integrate by parts on $D_1$ (using $\operatorname{div}v=0$ and homogeneous boundary conditions) to show
	\[
	-\frac{1}{\varepsilon\|v\|_{L_2(D_1)}^\beta}\int_{D_1} v\cdot\tilde\Delta v\,dx
	= \frac{\|\nabla v\|_{L_2(D_1)}^2}{\varepsilon\|v\|_{L_2(D_1)}^\beta}.
	\]
	By the Poincare inequality on $D_1$ there is $C_P>0$ with $\|v\|_{L_2(D_1)}\le C_P\|\nabla v\|_{L_2(D_1)}$, hence
	\[
	\frac{\|\nabla v\|_{L_2(D_1)}^2}{\varepsilon\|v\|_{L_2(D_1)}^\beta}
	\ge \frac{1}{\varepsilon C_P^\beta}\|\nabla v\|_{L_2(D_1)}^{2-\beta}.
	\]
	Collecting estimates and ensuring $\theta_2+\theta_3+\theta_4<1$ yields
	\begin{equation}\label{St 19}
		{\mu}\|\tilde\Delta v\|_{L_2(D)}^2 + \frac{1}{\varepsilon}\|\nabla v\|_{L_2(D_1)}^{2-\beta}
		\le C\big(\|v_t\|_{L_2(D)}^2 + \|\nabla v\|_{L_2(D)}^6 + \|f\|_{L_2(D)}^2\big),
	\end{equation}
	where $C$ depends on the choice of the small parameters and on $C_{\rm reg}$, $C_P$, etc.
	
Adding inequality \eqref{St 19} to \eqref{St 17} and dividing by $\theta_1>0$, we obtain  
\[
\| v^{\varepsilon}_{t} \|^{2}_{L_2(D)} + \frac{1}{\varepsilon} \|\nabla v^{\varepsilon}\|^{2-\beta}_{L_2(D_1)} + \frac{d}{dt}\Big( {\mu}\|\nabla v^{\varepsilon} \|^{2}_{L_2(D)} + \tfrac{1}{\varepsilon}\|v^{\varepsilon}\|^{2-\beta}_{L_2 (D_1)}\Big) 
\leq C\big( \|\nabla v^{\varepsilon}\|^{6}_{L_2 (D)} + \|f\|^{2}_{L_2 (\Omega)}\big).
\]
Setting  
\[
y(t)= {\mu}\|\nabla v^{\varepsilon} \|^{2}_{L_2(D)} + \tfrac{1}{\varepsilon}\|v^{\varepsilon}\|^{2-\beta}_{L_2 (D_1)},
\]
we get
\[
\frac{dy}{dt}\leq C \left(y^3(t)+\|f\|^{2}_{L_2 (\Omega)} \right),
\]
{with $C$ depending also on $\mu$.}
Integrating over $(0,t)$ gives  
\[
y(t)-y(0)\leq C\int_0^t y^3(\tau)\,d\tau + C\|f\|^{2}_{L_2 (Q_T)}.
\]
Let  
\[
z(t)=\int_0^t y^3(\tau)\,d\tau, \qquad B=C\|f\|^{2}_{L_2(Q_T)}.
\]
Then $y(t)\le y(0)+Cz(t)+B$ and hence  
\[
z'(t)=y(t)^3 \le (y(0)+Cz(t)+B)^3.
\]
Defining $R(t)=y(0)+Cz(t)+B$, we obtain
\[
R'(t)=Cz'(t)\le C R^3(t), \qquad R(0)=y(0)+B.
\]
By comparison with the solution of $S'(t)=C S^3(t)$, $S(0)=R(0)$, we find
\[
R(t)\le \left[-2Ct+(y(0)+B)^{-2}\right]^{-1/2}.
\]
Thus, choosing $T_0>0$ such that $-2CT_0+(y(0)+B)^{-2}>0$, we deduce that $R(t)$, and therefore $y(t)$, remain bounded on $[0,T_0]$. This yields the estimate
\begin{equation} \label{St 20}
	\|\nabla v^{\varepsilon} \|^2_{L_\infty(0,T_0;L_2(D))} + \tfrac{1}{\varepsilon} \| v^{\varepsilon} \|^{2-\beta}_{L_\infty(0,T_0;L_2(D_1))} + \| v^{\varepsilon}_{t} \|^2_{L_2(0, T_0; L_2(D))} + \tfrac{1}{\varepsilon} \int_0^{T_0} \|\nabla v^{\varepsilon} \|^{2-\beta}_{L_2(D_1)}\,dt \leq C.
\end{equation}
Finally, from \eqref{St 19} it also follows that
\begin{equation} \label{St 21}
	\|\tilde{\Delta} v^{\varepsilon} \|^2_{L_2(0,T_0;L_2(D))} \leq C
\end{equation}
for some $T_0>0$
depending on the solution of the homogeneous equation \eqref{St 1}--\eqref{St 3}, where the value \( T_0 \) is determined by the given data of the original problem \eqref{St 1}--\eqref{St 3}. 

Thus, we obtain estimate \eqref{St 15}. The  further proof of the local existence of a strong solution to the auxiliary problem \eqref{St 8}--\eqref{St 10} % and its convergence to the strong solution of the original problem \eqref{St 1}--\eqref{St 3} as \(\varepsilon \to 0\) 
is carried out based on estimates  \eqref{St 15} similarly to \cite[{section~III.3}]{antontsev1989boundary} and does not present any particular difficulty.
\medskip

\textbf{Proof of uniqueness of the strong solution.} Assume there exist two strong solutions $(v_1^\varepsilon,P_1^\varepsilon)$ and $(v_2^\varepsilon,P_2^\varepsilon)$ and set
$\omega=v_1^\varepsilon-v_2^\varepsilon$, $\pi=P_1^\varepsilon-P_2^\varepsilon$. Then from \eqref{St 8}--\eqref{St 10} we get
\begin{equation}\label{St22corr}
	\omega_t+(\omega\cdot\nabla)v_1^\varepsilon+(v_2^\varepsilon\cdot\nabla)\omega=\mu\Delta\omega-\nabla\pi
	-\frac{\xi(x)}{\varepsilon}\Big(\frac{v_1^\varepsilon}{\|v_1^\varepsilon\|_{L_2(D_1)}^\beta}-\frac{v_2^\varepsilon}{\|v_2^\varepsilon\|_{L_2(D_1)}^\beta}\Big),
\end{equation}
with $\operatorname{div}\omega=0$, $\omega|_{t=0}=0$, $\omega|_{S_1}=0$. Multiply \eqref{St22corr} by $\omega$ and integrate over $D$ to obtain
\[
\frac12\frac{d}{dt}\|\omega\|_{L_2(D)}^2+\mu\|\nabla\omega\|_{L_2(D)}^2
+\frac{1}{\varepsilon}\int_{D_1}\Big(\frac{v_1^\varepsilon}{\|v_1^\varepsilon\|^\beta}-\frac{v_2^\varepsilon}{\|v_2^\varepsilon\|^\beta}\Big)\!\cdot\omega\,dx
=-\int_D(\omega\cdot\nabla)v_1^\varepsilon\cdot\omega\,dx.
\]
The convection term is estimated by the Gagliardo--Nirenberg and Young inequalities:
\[
\Big|\int_D(\omega\cdot\nabla)v_1^\varepsilon\cdot\omega\,dx\Big|
\le C\|\nabla v_1^\varepsilon\|_{L_2}\|\omega\|_{L_2}^{1/2}\|\nabla\omega\|_{L_2}^{3/2}
\le\frac\mu2\|\nabla\omega\|_{L_2}^2 + C\|\omega\|_{L_2}^2.
\]
For the damping term set $a=\|v_1^\varepsilon\|_{L_2(D_1)}$, $b=\|v_2^\varepsilon\|_{L_2(D_1)}$. Using Cauchy--Schwarz,
$\int_{D_1}v_1^\varepsilon\cdot v_2^\varepsilon\,dx\le ab$, we obtain the monotonicity estimate
\[
\frac{1}{\varepsilon}\int_{D_1}\Big(\frac{v_1^\varepsilon}{a^\beta}-\frac{v_2^\varepsilon}{b^\beta}\Big)\cdot\omega\,dx
\ge\frac{1}{\varepsilon}(a-b)(a^{1-\beta}-b^{1-\beta})\ge0.
\]
Combining the estimates yields
\[
\frac{d}{dt}\|\omega\|_{L_2(D)}^2 \le C\|\omega\|_{L_2(D)}^2,
\]
and since $\omega(0)=0$ the Gronwall lemma gives $\omega\equiv0$ on $[0,T]$. Finally, from \eqref{St22corr} and $\omega\equiv0$ we infer $\nabla\pi=0$ (indeed the right-hand side of \eqref{St22corr} vanishes in $W^{-1}_2$), so the pressure difference is constant. Hence the strong solution is unique.
   Thus, the theorem is proved.
\end{proof}

\begin{thebibliography}{99}
\bibitem[antont]{antontsev1989boundary} {\sc S.~N. Antontsev, A.~V. Kazhiktov, and V.~N. Monakhov}, {\em Boundary value problems in mechanics of nonhomogeneous fluids}, vol.~22, Elsevier, 1989.
\end{thebibliography}
\end{document}
